\section{Datos y cómputo: cuándo supera ViT a la línea base convolucional}

El diagrama de la arquitectura no basta para responder si ViT es eficaz. A continuación, la clase compara distintos conjuntos de datos de preentrenamiento y presupuestos de cómputo, y muestra que ViT puede quedar por debajo de una CNN con regularización fuerte cuando hay pocos datos, pero logra mejor transfer con más datos y una recipe adecuada. Aquí el objeto causal es el sistema de entrenamiento completo, no el componente de attention por sí solo.

\subsection{Scaling with data: los supuestos a priori débiles necesitan más evidencia}

Cuando el preentrenamiento se amplía de ImageNet a ImageNet-21k y JFT-300M, la mejora de rendimiento de ViT es más notoria. De forma intuitiva, el fuerte supuesto a priori convolucional aporta sample efficiency en el régimen de pocos datos; al crecer la escala de los datos, un modelo más flexible puede aprender estructuras más adecuadas a la distribución de la tarea, y el costo de la falta de supuestos a priori se compensa poco a poco con los datos.

\lecturefigure{slide-19-scaling-with-data.jpg}{Resultados de ViT y BiT en ImageNet y en linear few-shot con distintas escalas de datos de preentrenamiento.}{Video oficial 30:24; Dosovitskiy et al. 2020.}

\begin{knowledgebox}{Lectura de la figura: primero el eje horizontal, luego el punto de cruce}
En la gráfica izquierda, el eje horizontal es la escala del conjunto de datos de preentrenamiento y el vertical es el top-1 de ImageNet; en la derecha, el eje horizontal es el número de muestras de JFT y el vertical es la linear few-shot accuracy. En el extremo de pocos datos, BiT/ResNet es más estable; en el extremo de muchos datos, la curva de ViT sube más rápido. Lo importante no es qué punto va adelante, sino que cada family de modelos responde a la escala de datos con una pendiente distinta. Las curvas no demuestran que con datos ilimitados se mantenga siempre la misma tendencia.
\end{knowledgebox}

\begin{warningbox}{La escala de datos no es un solo número de «cantidad de imágenes»}
La escala efectiva también depende de label noise, class diversity, long-tail coverage, duplicate rate, resolution, domain breadth y augmentation. 300M muestras web muy repetidas no equivalen a 300M unidades de información independientes. Al comparar scaling curves conviene informar las estrategias de limpieza y muestreo.
\end{warningbox}

\subsection{Frontera de preentrenamiento: los resultados de un artículo son producto conjunto del sistema}

La posición de un modelo en la performance frontier histórica depende en conjunto del dataset, el hardware, el optimizer, la training duration y la hyperparameter search. Que el punto de ViT quede por encima en la gráfica de evolución de modelos de ImageNet indica que esa recipe es competitiva; pero no es una ablación «architecture-only», ni implica que cualquier equipo pueda reproducir la misma ganancia con un presupuesto menor.

\lecturefigure{slide-20-vit-pretraining.jpg}{ViT forma un nuevo punto de frontier en la gráfica histórica de ImageNet accuracy contra escala de preentrenamiento.}{Video oficial 30:32.}

\begin{knowledgebox}{Lectura de la figura: la frontier solo compara las propuestas graficadas}
En la gráfica de dispersión, «más arriba» significa mayor accuracy y «más a la derecha» suele significar más datos o mayor escala. La frontier indica qué es mejor entre las propuestas reportadas, no que los ejes de Pareto estén completos; si no se marcan energy, latency, memory, annotation cost y tuning attempts, la gráfica no puede responder sobre eficiencia de despliegue ni costo de investigación.
\end{knowledgebox}

\subsection{Position embedding: ¿la estructura espacial se aprende o se escribe de antemano?}

Sin position embedding, la self-attention es aproximadamente permutation-equivariant respecto al orden de los tokens y no puede distinguir el orden espacial de un mismo conjunto de patches. ViT aprende un conjunto de vectores de posición absoluta; visualizar su similitud permite comprobar si el modelo recupera de la imagen bidimensional las relaciones de vecindad y la estructura de cuadrícula.

\lecturefigure{slide-21-position-embeddings.jpg}{La similitud de las codificaciones de posición aprendidas muestra una estructura de vecindad bidimensional, y se compara con distintas formas de inicialización/compartición.}{Video oficial 33:56.}

\begin{knowledgebox}{Lectura de la figura: el mapa de calor no prueba que «el modelo entiende el espacio»}
Si la embedding similarity de cada posición con las demás es localmente suave, el entrenamiento recuperó la estructura de vecindad bidimensional. La tabla de la derecha compara el rendimiento de distintos diseños de position embedding. Esta evidencia solo indica que la representación contiene regularidades espaciales; no prueba que el modelo haya adquirido geometría a nivel de objeto, relaciones causales ni una scene understanding completa; para ello se necesita el respaldo de attention patterns, probes y downstream behavior.
\end{knowledgebox}

\teachervoice{En la sesión de Q\&A, Lucas vuelve a esta slide y explica que la función principal de la position embedding es distinguir posiciones, y que la estructura de adyacencia local puede surgir durante el entrenamiento. Nota de clase: la visualización es un indicio del mecanismo, no una prueba causal; debe combinarse con experimentos de eliminación, permutación, interpolación y transferencia de resolución.}

\subsection{Compute-normalized comparison}

Si solo se compara por número de parámetros, el cómputo real de ViT, Mixer y CNN puede diferir mucho. La clase observa el transfer performance sobre el eje de compute: con FLOPs similares, la estructura del modelo y la recipe de entrenamiento determinan si el cómputo se convierte en representaciones transferibles. Antes de leer la gráfica hay que separar la cantidad teórica de operaciones del costo real del sistema: los mismos FLOPs pueden corresponder a formas de matriz, memory traffic, grado de paralelismo y comunicación completamente distintos, por lo que aquí solo se compara la compute efficiency a nivel algorítmico.

\lecturefigure{slide-22-scaling-with-compute.jpg}{Cada modelo forma una frontera de eficiencia distinta entre el compute de preentrenamiento y la transfer accuracy.}{Video oficial 39:24.}

\begin{knowledgebox}{Lectura de la figura: los FLOPs no son tiempo ni costo}
El eje horizontal de approximate compute facilita la comparación algorítmica, pero los mismos FLOPs pueden tener distinta latency, memory traffic y utilization en hardware diferente. Attention, convolution y MLP tienen distinta kernel efficiency; además, la comunicación distribuida no necesariamente se contabiliza. La gráfica puede responder «quién es más preciso con un presupuesto de cómputo teórico», pero no responde directamente «qué sistema es más barato o más rápido».
\end{knowledgebox}

\subsection{Resumen de la sección}

La ventaja de ViT aparece en el regime donde los datos, la regularization y el compute son suficientes. Las curvas de scaling de datos muestran la pendiente de respuesta, la gráfica de frontier muestra la competitividad histórica, la position embedding aporta indicios de estructura espacial y la gráfica de compute recuerda que el número de parámetros no es la única medida de costo.
